\section{Vecindad: utilidad dinámica por vecindad}
\label{sec:exp_e1}
%%%%%%%%%%%%%%%%%%%%%%%%%%%%%%%%%%%%%%%%%%%%%%%%%%%%%%%%%%%%%%%%%%%%%%%%%%%%%%%%

\textcolor{red}{%
[SCAFFOLDING --- Esta sección presenta los experimentos de validación para la estrategia Vecindad, que calcula el costo dinámico $D(v)$ contando el número de personas detectadas dentro de un radio fijo $r$ alrededor del centroide de la frontera $v$:
\[
    D_{\text{Vec}}(v) = \bigl|\{ p_i : \| c_v - p_i \|_2 \leq r \}\bigr|
\]
Los casos de prueba son TC001–TC010, con tres réplicas por caso (V001–V030). Se comparan los resultados con la línea de base Greedy (G001–G010) en las mismas 10 configuraciones de entorno y densidad de actores.
Los hallazgos principales son: Vecindad supera a Greedy en celdas fantasma en 6/10 configuraciones, en entropía en 5/10 y en frames en 3/10; por puntuación compuesta, Vecindad es competitiva o mejor en 5/10 configuraciones.
--- Completar con análisis detallado, tablas y figuras de Vecindad vs Greedy.%
]}

%%%%%%%%%%%%%%%%%%%%%%%%%%%%%%%%%%%%%%%%%%%%%%%%%%%%%%%%%%%%%%%%%%%%%%%%%%%%%%%%
